\section{Generator, constitution, execution, environment, and verifier}
\label{sec:deep-developmental}

A developmental system should not use one mechanism as imagination, judge, actuator, world, and scorekeeper. The minimum separation is
\begin{equation}
GENERATOR\neq CONSTITUTION\neq EXECUTOR\neq ENVIRONMENT\neq VERIFIER.
\label{eq:deep-separation}
\end{equation}
The generator proposes possibilities. The constitution evaluates admissibility, authority, and protected constraints. The executor realizes an authorized local transition. The environment includes dynamics and other independently authoring centers. The verifier compares returned consequence with claims and updates the model/policy boundary.

This separation does not require five software modules. It requires that interventions on one function need not silently rewrite the criteria of the others.

\paragraph{Critic-generator collapse.}
When the same process generates candidates and defines the only admissible evaluation vocabulary, novelty can shrink to whatever the critic already knows how to recognize. Conversely, an unconstrained generator can overwhelm verification with cheap novelty. A developmental architecture therefore needs both generative frontier and corrective frontier.

\paragraph{Founder removal as maturation.}
Let $H_f$ denote founder intervention. For a mature artifact, useful capability should remain above a task threshold as $H_f$ decreases, while correction capacity and provenance fidelity do not collapse:
\begin{equation}
H_f\downarrow,\quad Performance\ge\theta_P,\quad CorrectionAccess\ge\theta_C.
\label{eq:deep-founder}
\end{equation}
A system that performs only while the founder continuously supplies hidden context has not generalized the function; it has merely externalized founder labor.

\section{Developmental-agent architecture}

One prospective engineering target is a provenance-bearing research system whose useful capability can improve while founder maintenance and reconstruction burden fall. A minimal architecture separates:
\begin{equation}
\begin{aligned}
&\text{semantic transformation}+\text{persistent provenance}\nonumber\\
&+\text{candidate generation}+\text{warrant/governance}\nonumber\\
&+\text{tools/environment action}+\text{return capture}\nonumber\\
&+\text{verification}+\text{bounded write-back}.
\end{aligned}
\end{equation}
External memory, adaptation, substrate plasticity, recursive self-modification, and AGI remain separate claims. Current engineering work does not collapse them.

A useful improvement ratio is
\begin{equation}
\mathcal{R}=\frac{\text{validated durable work gain}}
{\text{maintenance}+\text{reconstruction}+\text{correction}+\text{intervention}}.
\end{equation}
Higher throughput with a faster-growing denominator is not the intended trajectory.

\section{Continuity, discontinuity, and carrier succession}
\label{sec:deep-continuity}

Continuity and discontinuity are complementary transformation conditions rather than opposites. Pure continuity can freeze a representation; pure discontinuity can erase every reconstructible invariant. A carrier succession succeeds when enough structure survives for later reconstruction while enough freedom remains for the successor to differ.

For consequence $\Delta$ transported through carriers $C_0,\ldots,C_n$, let $I_k$ be an invariant set retained across $C_k\rightarrow C_{k+1}$ and $N_k$ newly introduced structure. Mature succession requires neither $I_k=\varnothing$ nor $N_k=\varnothing$ by default. The research problem is to find the smallest invariant sufficient for regenerative reconstruction under later transformations.

This is the CBRD-side form of the companion paper's claim that carrier exit does not imply consequence deletion. It also supplies a concrete falsifier: if every useful continuation requires preservation of the original carrier's exact representation and authority, the alleged succession has not occurred.

\section{Evidence topology and current research boundary}
\label{sec:evidence-current}

A claim record is incomplete when it states only content and verdict. At minimum, record:
\begin{equation}
\begin{aligned}
\langle &content,register,status,source,scope,assumptions,exposure,\\
&alternatives,discriminator,revision\ condition,dependency\ closure,\\
&prohibited\ updates\rangle.
\end{aligned}
\end{equation}

Evidence classes include phenomenological, participant/naturalistic, historical, formal, engineering, synthetic, external empirical, and theological routes. Route existence does not transfer warrant between classes.

Hard negative closure for an executable result requires:
\begin{equation}
\begin{aligned}
claim/preregistration
&\rightarrow executable\ bytes
\rightarrow environment/runtime
\rightarrow raw\ outputs\\
&\rightarrow analysis/aggregation
\rightarrow result
\rightarrow scoped\ interpretation.
\end{aligned}
\label{eq:negative-proof-chain}
\end{equation}
A markdown verdict, frozen label, or successful build is evidence only for what it actually records.

The current accepted research boundary is:
\begin{itemize}
\item S0 frozen target and S1 comparison constitution;
\item S2--S4 accepted at their exact scopes;
\item corrupted S5 invalid and non-propagating;
\item clean Integrated S5--S7 accepted;
\item final S8--S9 accepted with audit narrowing;
\item S10 accepted with narrowing;
\item S11 executed, restarted, frozen, and participant-accepted with narrowing;
\item post-S11 recanonicalization is the current live successor state.
\end{itemize}

Current important negative/evidence statuses include:
\begin{itemize}
\item P5 inverse-neighbor/canonical-cycle privilege is a scoped negative in its reproducible tested representation, not a universal theorem;
\item historical HSF-v1 and HSF-v2 runs remain adverse synthetic evidence, while hard-negative authority is withheld at their present proof-chain gaps;
\item exact-six natural-kind/minimality remains \withhold;
\item conditional six-cycle algebra remains mathematically valid only under its assumptions and has no empirical privilege from syntax alone.
\end{itemize}

The post-S11 residual contribution is \emph{SUPPORTED-LOCAL / BOUNDED} at tested epistemic/governance/transport indices. Cross-World-entry generalization remains prospective. Historical uniqueness and universal formal irreducibility are not established.

\begin{readercheck}
For one result you regard as decisive, draw its entire evidence chain. Mark the first edge you cannot independently reconstruct. That edge limits the authority of the conclusion, even if the conclusion later turns out to be correct.
\end{readercheck}

\section{Prospective empirical program}

The current program is organized around discriminators rather than broad confirmation.

\subsection{Recursive sufficiency under future use}
Construct pairs of states that are currently compressible but diverge after a model-guided action reactivates source, history, or exposure. Compare with state-abstraction, bisimulation, predictive-state, and causal-abstraction baselines.

\subsection{Claim-relative non-preauthorship}
Construct coupled systems with controlled degrees of model influence on later observation. Test whether the proposed measure distinguishes genuinely corrective variation from self-conditioned surprise better than source count or raw disagreement.

\subsection{Compiler and RegisterBridge factorization}
Encode the same transformations using enriched contract/effect/lens or bidirectional-transformation formalisms. Compare composition, reversal, loss, authorization separation, evidence obligations, and failure propagation. If behavior is preserved, retain project language as a schema layer rather than claiming a new underlying calculus.

\subsection{AntiGrammar transfer}
Recover a transformation invariant from one representation and test held-out transfer and causal control against strong probing, steering, geometry, and metamorphic baselines.

\subsection{Functional differentiation}
Blindly discover failure clusters under role fusion/ablation, then compare the recovered decomposition against the historical six-role basis and alternatives. No-clean-mapping is an admissible result.

\subsection{Regenerative transport and founder removal}
Remove vocabulary, founder rescue, or original carrier access while measuring reconstruction, criticism, mutation, provenance fidelity, novel transfer, and correction quality.

\section{Formal frontier}

The highest-leverage open problems are currently:
\begin{enumerate}
\item first consequential differentiation $\delta$ without presupposed cut;
\item exact typed claim/warrant/status calculus;
\item reusable negative-result semantics;
\item recursive sufficiency under self-conditioning and policy change;
\item claim-relative causal non-preauthorship;
\item Compiler composition, partiality, loss, obligation, and failure calculus;
\item RegisterBridge composition/reversal and comparison with enriched bidirectional formalisms;
\item AntiGrammar identification and transfer theory;
\item exact Perspectivalization calculus and Observer reconstruction;
\item general projection/section/reconstruction theory, with the P-versus-NP program only one specialization.
\end{enumerate}

A result may be that some project-specific vocabulary is unnecessary. That counts as progress.

\section{Reduction and retirement}

CBRD should shrink whenever established language preserves the relevant consequence more cleanly.
\begin{enumerate}
\item Retire a provenance coordinate if source-swap and future-use tests show it never matters for the declared task.
\item Retire a named failure construct if ordinary mediation, routing, control, or causal graphs diagnose the same intervention point.
\item Retire AntiGrammar if existing probing/steering/metamorphic methods preserve its transfer and control value.
\item Treat RegisterBridge and Compiler as schema vocabulary if enriched established formalisms preserve all decision behavior.
\item Merge functional roles when deletion/fusion produces no distinctive failure.
\item Keep regenerative-transport language outside empirical claims if it cannot be operationalized without founder judgment.
\item Withhold hard-negative closure when executable evidence chains are incomplete.
\end{enumerate}

Global failure occurs if the framework cannot become smaller, if every counterexample is redescribed as support, if status labels outrun their typed objects, or if its own prose safeguards substitute for operative controls.

\section{Conclusion}

CBRD studies a recurring technical problem: a representation can become part of the system that later appears to verify it. That makes provenance, projection, authorization, and future-use sufficiency more important, not less. The framework's strongest claims are therefore separations and tests rather than cosmic conclusions.

The portable core is: localize where exclusion entered; preserve provenance when it changes future consequence; judge compression under the transformations induced by using it; distinguish returned consequence from self-authored confirmation; type claims before verdicts; keep negative results scoped to their evidence chain; separate compilation from authorization and verification; treat projection as information loss unless a lawful reconstruction is earned; and delete project vocabulary when mature alternatives do the same job better.

\begin{readercheck}
State three experiments that could make CBRD materially smaller without using the terms \emph{recursive sufficiency}, \emph{RegisterBridge}, \emph{AntiGrammar}, or \emph{Compiler}. If the framework cannot describe its own reduction without those names, it is not yet portable.
\end{readercheck}
